\section{Theta transformation and the functional equation}
\label{sec:theta-functional}

This section uses the gamma gateway proved in the preceding material:
Euler's integral on a half-plane with a supplied positive margin,
$\Gamma(z+1)=z\Gamma(z)$, $\Gamma(1/2)=\sqrt\pi$, reflection,
and duplication.  These are dependencies on human-readable finite
proofs, not assertions that the corresponding Lean development has
already been checked.  All integrals below denote prescribed rational
quadrature algorithms with the displayed endpoint estimates.

\subsection{An effective Fourier uniqueness lemma}

For an integer $N\ge1$, define the finite trigonometric polynomial
\[
 K_N(u)=\frac1N\left|\sum_{j=0}^{N-1}e^{2\pi iju}\right|^2
       =\sum_{|j|<N}\left(1-\frac{|j|}N\right)e^{2\pi iju}.
\]
The second equality follows by multiplying two finite sums and counting
pairs with difference $j$.  Consequently $K_N\ge0$ and
$\int_0^1K_N(u)\,du=1$: integrate each exponential, using
$e^{2\pi ij}=1$ and its elementary antiderivative when $j\ne0$.
For $0<\delta\le1/2$ and
$\operatorname{dist}(u,\mathbb Z)\ge\delta$, finite geometric summation
and $\sin(\pi v)\ge2v$ on $[0,1/2]$ give
\[
 K_N(u)=\frac{\sin^2(\pi Nu)}{N\sin^2(\pi u)}
             \le\frac1{4N\delta^2}.                 \tag{TF1}
\]
The sine inequality is the finite second-derivative comparison with the
chord from $0$ to $\pi/2$; the trigonometric gateway proves it by rational
meshes.  No exact comparison with a zero of sine occurs in (TF1), since
its denominator has the displayed positive lower bound.

Let $f$ be a one-periodic function with $|f|\le M$ and Lipschitz
constant $L$ for circular distance, both rational bounds.  Splitting
its convolution with $K_N$ at circular distance $\delta$ proves
\[
 \left|\int_0^1K_N(u)f(x-u)\,du-f(x)\right|
       \le L\delta+\frac{M}{2N\delta^2}.             \tag{TF2}
\]
On the near part use $L\delta$ and mass at most one; on the remaining
part use $2M$, (TF1), and length at most one.  Thus, to obtain error
$2^{-q}$, it suffices to take
\[
 \delta=\frac{2^{-q-1}}{L+1},\qquad
 N\ge M(L+1)^2 2^{3q+2}.                             \tag{TF3}
\]
Integrating the finite expression for $K_N$ shows that the convolution
is the finite weighted Fourier sum
\[
 \sum_{|j|<N}\left(1-\frac{|j|}N\right)\widehat f(j)e^{2\pi ijx},
 \qquad \widehat f(j)=\int_0^1f(u)e^{-2\pi iju}\,du.
                                                               \tag{TF4}
\]
Equations (TF1)--(TF4), including each integral error, are finite
certificates.  They replace any appeal to a general Fourier-series
convergence theorem.  In particular two such functions with the same
Fourier coefficients are equal, by evaluating (TF2) at an arbitrary
precision.

\subsection{The Gaussian transform without complex contour integration}

Fix $t$ with $1/A\le t\le A$, where $A\ge1$ is an integer.  Put
\[
 G_t(b)=\int_{-\infty}^{\infty}
              e^{-\pi t x^2}e^{-2\pi ibx}\,dx \qquad(b\in\mathbb R).
\]
For every integer $m\ge0$ and $T\ge1$,
\[
 \int_{|x|\ge T}|x|^m e^{-\pi t x^2}\,dx
       \le 2m!(2A)^{m+1}e^{-T/(2A)}.                \tag{TG1}
\]
Indeed $\pi t x^2\ge |x|/A$ there, and
$|x|^m\le m!(2A)^m e^{|x|/(2A)}$.  Integrating the resulting
exponential on each half-line gives (TG1).  The extra constant is
harmless.  It provides primitive recursive cutoffs for differentiation
in $b$ and for every integration by parts used below.

Positive-real substitution in the gamma integral, together with
$\Gamma(1/2)=\sqrt\pi$, gives $G_t(0)=t^{-1/2}$.  Differentiate the
finite integral and then apply (TG1) with $m=1,2$ to its remainder.
Integration by parts in $x$ gives
\[
 G_t'(b)=-\frac{2\pi b}{t}G_t(b).
\]
For clarity, if $f(x)=e^{-\pi t x^2}$, then
$f'(x)=-2\pi txf(x)$ and
$\int f'(x)e^{-2\pi ibx}\,dx=2\pi ibG_t(b)$;
the boundary is bounded by (TG1).  Differentiating
$e^{\pi b^2/t}G_t(b)$ therefore gives zero.  The finite fundamental
theorem of calculus on the bounded segment from $0$ to $b$ proves
\[
 \boxed{G_t(b)=t^{-1/2}e^{-\pi b^2/t}.}             \tag{TG2}
\]
All frequency and exponential factors have rational bounds once a
bound for $|b|$ is supplied.  Thus this argument contains an algorithm
for every precision; it does not invoke an existence theorem for ODEs.

\subsection{Poisson summation for the one Gaussian that is needed}

Consider the one-periodic function
\[
 P_t(x)=\sum_{n\in\mathbb Z}e^{-\pi t(x+n)^2}.
\]
Uniformly for $0\le x\le1$, the omitted terms with $|n|>N$, $N\ge2$,
have total at most
\[
 8A e^{-N/(2A)}.                                    \tag{TG3}
\]
For such $n$, $|x+n|\ge |n|/2$, so a geometric series with ratio
$e^{-1/(2A)}$ suffices; use
$1-e^{-1/(2A)}\ge1/(4A)$.  The same estimate with a factor $2|n|$
proves convergence of the differentiated series.  In particular the
safe global bounds
\[
 |P_t(x)|\le16A,\qquad |P_t'(x)|\le1024A^3           \tag{TG4}
\]
follow from
$\sum_{n\ge1}nq^n=q/(1-q)^2$ and $\pi<4$.
These identities for the geometric series follow first for finite sums,
with their explicit geometric remainders.

Integrate a finite periodization against $e^{-2\pi ikx}$, substitute
$u=x+n$ in its $n$th term, and use $e^{2\pi ikn}=1$.  Its intervals
join to $[-N,N+1]$.  Estimates (TG1) and (TG3) therefore identify the
Fourier coefficient with (TG2):
\[
 \widehat P_t(k)=t^{-1/2}e^{-\pi k^2/t}.
\]
The coefficients have the useful bound
\[
 \sum_{k\in\mathbb Z}|k|\,|\widehat P_t(k)|\le8A^3. \tag{TG5}
\]
To check it, use $t^{-1/2}\le A$,
$e^{-\pi k^2/t}\le e^{-|k|/A}$ for $|k|\ge1$, and
$1-e^{-1/A}\ge1/(2A)$.  The difference between the full coefficient
sum and its Fej\'er weighted sum is consequently at most $8A^3/N$:
for $|k|\ge N$ the missing weight $1$ is still at most $|k|/N$.
Combining this bound with (TF2)--(TF3), at $x=0$, proves
\begin{equation}
 \boxed{\displaystyle
 \theta(t):=\sum_{n\in\mathbb Z}e^{-\pi n^2t}
       =t^{-1/2}\theta(1/t)\qquad(t>0).}
 \label{eq:theta-transformation}
\end{equation}
This is the required Poisson summation formula.  Its proof used only
finite Fourier polynomials, one explicitly normalized Gaussian
integral, and specified error bounds.

\subsection{An entire integral and the completed zeta function}

Write $\psi(t)=(\theta(t)-1)/2=\sum_{n\ge1}e^{-\pi n^2t}$.
For $t\ge1$, elementary geometric summation gives
\[
 0<\psi(t)\le2e^{-3t},\qquad
 \sum_{n>N}e^{-\pi n^2t}\le2e^{-3(N+1)t}.            \tag{TH1}
\]
Define
\begin{equation}
 A(s)=\int_1^\infty\psi(t)
             \bigl(t^{s/2}+t^{(1-s)/2}\bigr)\,\frac{dt}{t}.
 \label{eq:theta-entire-integral}
\end{equation}
For an integer $R\ge1$, $|s|\le R$, and every derivative order $j$,
its tail after $T\ge1$ is bounded by
\[
 \int_T^\infty\left|
 \frac{\partial^j}{\partial s^j}
 \left[\psi(t)(t^{s/2}+t^{(1-s)/2})/t\right]\right|dt
       \le2(R+j)!e^{-2T}.                           \tag{TH2}
\]
Indeed the integrand is at most
$4t^{R+j}e^{-3t}$, since $\log t\le t$, and
$t^{R+j}\le(R+j)!e^t$.  On $[1,T]$ use the finite theta sum and a
uniform mesh.  A deliberately generous Lipschitz bound for the
$j$th integrand is
\[
 C_{R,j,T}=2^{R+j+10}(R+j+2)!(T+1)^{R+j+4}.          \tag{TH3}
\]
It follows by the product rule, $|\psi'|\le8$, and the same geometric
majorants; it also bounds the finite theta approximants.
For precision $p$, one can take
$T=p+2(R+j)!+R+j+10$, take the theta cutoff sufficiently large that
$4(T+1)^{R+j+2}e^{-3N}\le2^{-p-6}$, and use at least
$2^{p+6}C_{R,j,T}(T+1)^2$ mesh intervals.  The theta cutoff can be
chosen explicitly as
$N=p+(R+j+4)(T+1)+10$; no unbounded search is needed.  Elementary
function evaluations are made to error at most
$2^{-p-6}/(T+1)$ per integrand value.

The integral has an everywhere convergent Taylor representation.
At a center $s_0$ its $j$th Taylor coefficient is
\[
 \frac1{2^j j!}\int_1^\infty\psi(t)
 \bigl(t^{s_0/2}+(-1)^jt^{(1-s_0)/2}\bigr)
          (\log t)^j\,\frac{dt}{t}.                 \tag{TH4}
\]
For any prescribed radius $r>0$, sum the absolute integrands with
weights $(2r)^j$.  The result is bounded by the same integrand with
an extra factor $t^r$, hence by (TH2) with a larger integer $R$.
If $C$ is that explicit integral majorant, the Taylor tail for
$|s-s_0|\le r$ after degree $J$ is at most $C2^{-J}$.
Expanding a finite exponential polynomial and adding this tail proves
(TH4) and termwise differentiation.  This proves the effective
entireness assertion, with a modulus, without a general theorem about
locally uniform limits.

For $\Re s>1$ the gamma substitution gives
\[
 \int_0^\infty\psi(t)t^{s/2-1}\,dt
        =\pi^{-s/2}\Gamma(s/2)\zeta(s).             \tag{TH5}
\]
Here the interchange is justified before taking any limit: after
summing $n\le N$, the absolute integral of the omitted sum is at most
\[
 \pi^{-\sigma/2}\Gamma(\sigma/2)
                       \frac{N^{1-\sigma}}{\sigma-1},
 \qquad \sigma=\Re s>1.                            \tag{TH6}
\]
This follows by the positive-real substitution in each term and the
already proved Dirichlet tail estimate.  Gamma bounds and a supplied
margin $\sigma-1\ge1/A$ make (TH6) an explicit primitive recursive
cutoff.

Split (TH5) at $1$.  The transformation
\eqref{eq:theta-transformation} says
\[
 \psi(t)=t^{-1/2}\psi(1/t)+(t^{-1/2}-1)/2.
\]
In the first term substitute $u=1/t$; integrate the two remaining
powers by their elementary antiderivatives.  Endpoint removal is
justified by (TH6) and the supplied margin.  The result is
\begin{equation}
 \pi^{-s/2}\Gamma(s/2)\zeta(s)
       =\frac1{s(s-1)}+A(s)\qquad(\Re s>1).
 \label{eq:theta-completion-initial}
\end{equation}
Define the meromorphic completion and its pole-free form by
\begin{align}
 \Lambda_\theta(s)&=\frac1{s(s-1)}+A(s),\label{eq:theta-lambda}\\
 \xi(s)&=\frac12+\frac12s(s-1)A(s).\label{eq:theta-xi}
\end{align}
The first formula is evaluated only when both denominators have
positive rational separation; the second is an entire name on every
bounded disk, including $0$ and $1$.

\subsection{Quantitative continuation, with a finite induction invariant}
\label{subsec:theta-overlap}

The equality on the initial half-plane must be identified with the
Euler--Maclaurin continuation already constructed.  Merely saying
``by the identity theorem'' would conceal the main weak-arithmetic
issue.  The following finite estimate supplies the needed replacement.

Suppose a function has convergent Taylor series on a disk of radius
slightly greater than $R$, is bounded by $H\le2^h$ on the disk of
radius $R$, and is bounded by $\varepsilon$ on the concentric disk of
radius $R/4$.  Integrating its Taylor polynomials around circles,
with their explicit tails, extracts each coefficient by finite Fourier
orthogonality.  Hence its $j$th coefficient has both bounds
$H/R^j$ and $\varepsilon(4/R)^j$.  Split its series after degree $J$.
On the disk of radius $R/2$ this proves
\begin{equation}
 |F|\le 2^{J+1}\varepsilon+2H\,2^{-(J+1)}.
 \label{eq:analytic-propagation}
\end{equation}
Circle integration here is just integration over $0\le u\le1$ of
$F(c+re^{2\pi iu})e^{-2\pi iju}$; the coefficient formula follows by
integrating finite power polynomials.  No Cauchy integral theorem or
choice of an unspecified zero is involved.

Take $T$ disks whose successive centers are separated by at most
$R/4$, with the same bound $H\le2^h$.  To achieve final error
$2^{-q}$, prescribe integers backwards by
\[
 a_T=q,\qquad a_j=2a_{j+1}+h+8,
 \qquad J_j+1=h+a_{j+1}+4.                          \tag{TH7}
\]
Then \eqref{eq:analytic-propagation} carries error $2^{-a_j}$ to
error at most $2^{-a_{j+1}}$.  Each next disk of radius $R/4$ lies
inside the previous disk of radius $R/2$.  Explicitly,
\[
 a_0=2^T(q+h+8)-h-8.
\]
The induction invariant is a bound on finite approximations at
specified precisions, not the quantified statement that all
coefficients vanish.  The two analytic errors sum to at most
$(3/16)2^{-a_{j+1}}$, leaving a reserve for numerical errors.
The outer coefficient bound gives
$|F'|\le16H/(9R)$ on the radius-$R/4$ disk. Consequently the coefficient
integrand
$F(c+(R/4)e^{2\pi iu})e^{-2\pi iju}$ has Lipschitz bound
$8(j+1)H$. A left mesh with at least $8(j+1)H2^b$ intervals gives
quadrature error at most $2^{-b}$. Take
$b=a_{j+1}+J_j+8$ and assign the same bound to function and circle-point
evaluation errors; their finite sum, including amplification by
$\sum_{j\le J_j}2^j<2^{J_j+1}$, fits inside the reserve.
All indices are primitive recursive. More explicitly, for a desired
endpoint approximation form the finite tree of circle sample points
needed by the successive coefficient extractions.  The mesh sizes and
Taylor degrees are fixed backwards by (TH7); a bounded scan of this
finite tree carries all local rational error bounds.  This is the
quantifier-free induction invariant used internally.  If $F$ is zero on the
initial disk at every prescribed precision, (TH7) makes it zero at
the endpoint at every prescribed precision.

Here is the application, including removal of possible poles.
Let $\zeta_{\rm EM}$ denote the Euler--Maclaurin construction, and
let $E(s)=(s-1)\zeta_{\rm EM}(s)$, evaluated at $1$ by its explicit
removable-pole formula.  For an integer $r\ge1$ put
\[
 P_r(s)=\prod_{j=0}^{r-1}(s/2+j),
\]
\[
 F_r(s)=s\pi^{-s/2}\Gamma(s/2+r)E(s)
            -P_r(s)\{1+s(s-1)A(s)\}.              \tag{TH8}
\]
Every factor is effectively analytic on the convex half-plane
$\Re s>-2r$.  For gamma this follows by its integral Taylor expansion;
for $E$ it follows by the Euler--Maclaurin Taylor bounds; for $A$ it
is (TH4).  The gamma recurrence and
\eqref{eq:theta-completion-initial} make $F_r=0$ on $\Re s>1$.

To reach a target $s$ with a rational certificate
$\Re s+2r\ge a>0$, join $2$ to $s$ by a straight segment.  Set
$R=\min(1,a)/4$ and choose
$T\ge4D/R$, where $D\ge|s-2|$ is a supplied rational bound.
All radius-$R$ disks stay in the half-plane, and the initial
radius-$R/4$ disk lies in $\Re s>1$.  A uniform rational bound $H$
for (TH8) is obtained directly from its displayed factors.  For
example, if $B\ge1$ is an integer with $|z|\le B$ on the disks and $|E(z)|\le E_B$, one may use
\[
 \begin{split}
 G&=\lceil4/a\rceil+1+2^{B+r+2}(B+r+1)!,\\
 H&=1+B4^BGE_B+(B+r)^r\{1+2B(B+1)B!\}.
 \end{split}                                                   \tag{TH9}
\]
Here $|A(z)|\le2B!$ follows from (TH2), and $E_B$ is the finite
Euler--Maclaurin bound obtained by taking its truncation point $1$
and order $m=B+2$.  In particular $E_B$ is computed by finite sums
of absolute Bernoulli coefficients, products bounded by
$(B+2m)^{2m}$, and an integral remainder bounded by the sum of those
coefficients.  One can take the integer $h=\lceil H\rceil$.
Formula (TH7) now proves $F_r(s)=0$ by an explicit finite error budget.

Consequently the two continuations agree.  On any chart where the
ordinary factors are defined this yields
\[
 \pi^{-s/2}\Gamma(s/2)\zeta_{\rm EM}(s)=\Lambda_\theta(s).
                                                               \tag{TH10}
\]
At gamma poles one uses (TH8), rather than dividing by a zero of
$P_r$.  This is a divisor-free statement valid across those points.
It is exactly the regularized meaning of (TH10).

\subsection{Functional equations and consequences}

Interchanging the two summands in \eqref{eq:theta-entire-integral}
gives $A(s)=A(1-s)$ at every precision.  The formulas above prove
\begin{equation}
 \boxed{\Lambda(s)=\Lambda(1-s),\qquad \xi(s)=\xi(1-s).}
 \label{eq:completed-functional-equation}
\end{equation}
Moreover
\[
 \xi(0)=\xi(1)=\tfrac12,\qquad
 \xi(\overline s)=\overline{\xi(s)}.
\]
The second identity follows term by term from real quadrature weights
and conjugation of exponentials.  Thus $\xi(1/2+it)$ is real for
real $t$, and every odd Taylor coefficient of $\xi$ at $1/2$ is zero.
Indeed substitute $s=1/2+h$ in the symmetry and compare coefficients
by the finite circle-extraction argument.  In derivative notation,
$\xi^{(2j+1)}(1/2)=0$ for all $j\ge0$.

On a closed substrip of $0<\Re s<1$, the gamma reflection and
 duplication identities are valid with supplied positive margins.
The completed equation, multiplied by $\Gamma(1-s/2)$, gives
\[
 \Gamma(s/2)\Gamma(1-s/2)=\frac\pi{\sin(\pi s/2)},
\]
\[
 \Gamma((1-s)/2)\Gamma(1-s/2)
                       =2^s\sqrt\pi\,\Gamma(1-s).
\]
Substitution proves
\begin{equation}
 \boxed{\displaystyle
 \zeta(s)=2^s\pi^{s-1}\sin(\pi s/2)
                        \Gamma(1-s)\zeta(1-s).}
 \label{eq:zeta-sine-functional}
\end{equation}
To extend this formula into $\Re s<1$, multiply its difference by
$s$.  The pole of $\zeta(1-s)$ at zero is then removed, and every
factor is analytic on this half-plane.  Apply exactly the finite
disk propagation (TH7), starting at $1/2$.  On arbitrary other bounded
charts use gamma recurrence to remove its finitely many possible
poles, and apply the same argument.  Thus
\eqref{eq:zeta-sine-functional} is a meromorphic identity, evaluated
without division at removable points.  Replacing $s$ by $1-s$ also
gives
\[
 \zeta(1-s)=2(2\pi)^{-s}\cos(\pi s/2)\Gamma(s)\zeta(s).
\]

For $m\ge1$, all factors on the right side of
\eqref{eq:zeta-sine-functional} are analytic near $s=-2m$, and the
sine factor vanishes.  Therefore $\zeta(-2m)=0$.  Differentiate the
finite Taylor identity once; only the derivative of sine survives:
\begin{equation}
 \boxed{\displaystyle
 \zeta'(-2m)=(-1)^m\frac{(2m)!}{2^{2m+1}\pi^{2m}}
                              \zeta(2m+1).}
 \label{eq:trivial-zero-derivative}
\end{equation}
Since $\zeta(2m+1)>1$ and $\pi<4$,
\[
 |\zeta'(-2m)|\ge2^{-6m-1}.                         \tag{TH11}
\]
This is an explicit apartness certificate.  Each negative even zero
is simple; the assertion is stronger than a bare negation that its
derivative equals zero.  On any compact part of $\Re s<0$ with a
positive separation from the negative even integers, the sine product,
the nonvanishing gamma product, and the right-half-plane lower bound
for $\zeta(1-s)$ similarly give a positive lower bound for $|\zeta(s)|$.
Consequently there are no further zeros on those apartness charts.
For real $x<0$ away from even integers, all factors other than sine in
\eqref{eq:zeta-sine-functional} are positive, so
$\operatorname{sign}\zeta(x)=\operatorname{sign}\sin(\pi x/2)$.

Finally let $\gamma$ denote Euler's constant.  The gamma product
and the Laurent construction already proved give, with effective
Taylor remainders,
\[
 \Gamma(1-s)=1+\gamma s+O(s^2),\qquad
 \zeta(1-s)=-1/s+\gamma+O(s).
\]
Their constant terms cancel in the product:
$\Gamma(1-s)\zeta(1-s)=-1/s+O(s)$.
Also
$2^s\pi^{s-1}=\pi^{-1}(1+s\log(2\pi)+O(s^2))$ and
$\sin(\pi s/2)=\pi s/2+O(s^3)$.  Multiply these finite
polynomials with their explicit remainders in
\eqref{eq:zeta-sine-functional}; at $s=0$ this gives
\[
 \boxed{\zeta(0)=-\tfrac12,\qquad
             \zeta'(0)=-\tfrac12\log(2\pi).}
\]
Here each $O(s^j)$ abbreviates a rational bound $C|s|^j$ on a fixed
small disk obtained from the cited Taylor majorants; it is not an
additional asymptotic assumption.

\subsection{Formalization interface for this section}

The Lean development should first formalize the finite lemmas:
weighted Fourier convolution (TF4), its error (TF2), the Gaussian tail
(TG1), periodization tail (TG3), and finite circle coefficient extraction.
The theta identity then becomes equality of two approximation programs
with the combined (TF2), (TG3), and (TG5) error budget.  Define $A$ by the
cutoffs following (TH3), define $\xi$ by
\eqref{eq:theta-xi}, and prove its Taylor coefficients using (TH4).
The most important reusable lemma is
\eqref{eq:analytic-propagation}, with its finite-list iteration (TH7).
It should be stated as a quantitative lemma for certified Taylor
charts, not imported as the classical analytic identity theorem.
Apply it to the explicit expression (TH8), then prove the symmetric
functional equation by the two-term interchange.  Gamma identities
supply the sine conversion, and finite Taylor coefficient arithmetic
supplies \eqref{eq:trivial-zero-derivative} and the values at zero.
These are precise proof obligations for a future full formalization;
this section does not claim an existing kernel-checked proof of them.
